%-------------------------------------------------------------------------------
%	SECTION TITLE
%-------------------------------------------------------------------------------

\cvsection{Опыт работы}

\begin{cventries}
\cventry
  {Java Backend Developer} % Job title
  {ПАО Сбербанк - команда СберОпрос} % Organization
  {Россия} % Location
  {Сентябрь 2023 — н.в. (2 Года)} % Date(s)
  {
    \begin{cvitems} % Description(s) of tasks/responsibilities
      \item {Разрабатываю бэкенд продукта для проведений качественных исследований через анкетирования - SberOpros, который собрал за 4 года 6500 анкет, 
      18 млн ответов и зарегистрировал 80 млн респондентов. Продукт позволяет внутренним сотрудникам Сбербанка создавать и проводить опросы на юридических лицах и сотрудниках 
      компании в формате ссылок, email-писем, а также виджета, встраиваемого в другие продукты}
      \item {Разработал интеграцию с сервисом для подбора респондентов, что повысило конверсию опросов.}
      \item {Реализовал сбор обратной связи от пользователей конструктора через виджет, что позволило улучшить качество обслуживания.}
      \item {Реализовал сбор пользовательских метрик (DAU/WAU/MAU, Cycle Time, LTR Виджета) и выгрузку отчётов, что улучшило пользовательский опыт и поддержку бизнеса.}
      \item {Провел онбординг 4 новых сотрудников и взял на себя обязанности по предрелизным делам на тестовых стендах.}
      \item {\textbf{Стек}: Java 17, Spring Boot 2.7, DDD, Clean Architecture, PostgreSQL, Spring Data JPA, Hibernate, REST, OpenAPI, Kafka, Keycloak, Kubernetes, Istio, Prometheus.}
    \end{cvitems}
  }
\end{cventries}




%-------------------------------------------------------------------------------
%	CONTENT
%-------------------------------------------------------------------------------
\cvsection{Проектная деятельность}

\begin{cventries}

%---------------------------------------------------------
\cventry
  {}
  {Vogorode}
  {\href{https://github.com/ArchieSW/vogorode}{\underline{Github}}, 2023}
  {}
  {
    \vspace{-5mm}
    \begin{cvitems}
      \item {Разработал систему управления заявками на услуги мастеров и садоводческих работ в микросервисной архитектуре.}
      \item {Организовал взаимодействие между сервисами, реализовал балансировку и отказоустойчивость.}
      \item {Интегрировал Kafka для асинхронного обмена сообщениями между сервисами.}
      \item {Настроил контейнеризацию и оркестрацию приложения.}
      \item {Внедрил мониторинг для повышения надежности и observability системы.}
      \item {Оптимизировал работу с базой данных, внедрив пагинацию, индексы и кастомные запросы через.}
      \item {\textbf{Стек}: Java 17, Spring Boot 2.7, PostgreSQL, Spring Data JPA, Hibernate, REST, gRPC, Prometheus, Grafana, Docker, Docker Compose}
    \end{cvitems}
  }

 \vspace{1em}

%---------------------------------------------------------
  \cventry
    {}
    {Распределенная обработка финансовых данных}
    {\href{https://github.com/HSECourseWork2022}{\underline{Github}}, 2023}
    {}
    {
      \vspace{-5mm}
      \begin{cvitems}
        \item {Командная курсовая работа по разработке etl системы обработки финансовых данных в реальном времени для аналитики.
         Занимался написанием микросервисов для скрапинга данных с api криптобирж (продьюсеры) и их обработки (консьюмеры).}
        \item {Написание микросервисов с использованием Spring Boot}
        \item {Связывание микросервисов при помощи Kafka}
        \item {\textbf{Стек}: Java 17, Spring Boot 2.7, PostgreSQL, Spring Data JPA / Hibernate, REST, Kafka, Docker, Docker Compose}
      \end{cvitems}
    }

%---------------------------------------------------------
\end{cventries}
